\definecolor{UnsafePurple}{HTML}{674ea7}
\newcommand{\unsafe}{\code{unsafe}\xspace}
\newcommand{\code}[1]{{\color{UnsafePurple}\texttt{\textscale{1}{#1}}}}
\newcommand{\qopt}{{$\star$}}
\newcommand{\internote}[1]{{\textbf{[}} \textit{#1} {\textbf{]}}}
\newcommand{\todonum}{{\color{red}[insert \#]}\xspace}
\newcommand{\tododate}{{\color{red}[insert date]}\xspace}
\newcommand{\tempdata}[1]{{\color{red}[#1]}\xspace}

\newcommand{\rsq}[2]{\paragraph{\textbf{RQ}\textbf{#1} - \textit{#2}}}
\newcommand{\CC}{C\texttt{++}}


\newcommand{\ilcite}[1]{(\hyperref[table:participants]{P#1})}
\NewDocumentEnvironment{pquote}{m}
{
\begin{quote}\itshape
}
{
\noindent \null \hfill {\normalfont \ilcite{#1} }
\end{quote}
}

\newcommand{\ilquote}[2]{``\textit{#1}''~\ilcite{#2}}

\newcommand{\datanote}[1]{{\textit{[Note]: #1}}}




\newcommand{\ArrayItemRounded}[1]{\FPeval{\thisitem}{round(clip(\ArrayItem{#1}),0)}\thisitem}


\newcommand{\rsqone}{\rsq{1}{What do developers state as their motivations for using Rust?}}
\newcommand{\rsqtwo}{\rsq{2}{What do Rust developers state as their motivations for using \unsafe code?}}
\newcommand{\rsqthree}{\rsq{3}{How do Rust developers reason about encapsulating \unsafe code?}}
\newcommand{\rsqfour}{\rsq{4}{How do Rust developers reason about memory safety across foreign boundaries?}}
\newcommand{\rsqfive}{\rsq{5}{How do Rust developers validate \unsafe code?}}
\newcommand{\rsqsix}{\rsq{6}{How is unsafe code perceived by the members of the Rust community who engage with it?}}

\usepackage[nomessages]{fp}% http://ctan.org/pkg/fp

\usepackage{xparse}
\usepackage{xfp}
\ExplSyntaxOn
\ior_new:N \l__lrnv_temp_ior
\str_new:N \l__lrnv_tmp_str
\seq_new:N \l__lrnv_key_seq
\seq_new:N \l__lrnv_value_seq
\prop_new:N \l__lrnv_array_prop
\cs_new_protected:Npn \__lrnv_csv_to_prop:n #1
  {
    \__lrnv_input_csv:n {#1}
    \seq_mapthread_function:NNN \l__lrnv_key_seq \l__lrnv_value_seq
      \__lrnv_set_prop:nn
  }
\cs_new_protected:Npn \__lrnv_set_prop:nn #1 #2
  { \prop_put:Nnn \l__lrnv_array_prop {#1} {#2} }
\cs_new_protected:Npn \__lrnv_input_csv:n #1
  {
    \ior_open:Nn \l__lrnv_temp_ior {#1}
    \__lrnv_read_csv_row_to_seq:NN \l__lrnv_temp_ior \l__lrnv_key_seq
    \__lrnv_read_csv_row_to_seq:NN \l__lrnv_temp_ior \l__lrnv_value_seq
    \ior_close:N \l__lrnv_temp_ior
  }
\cs_new_protected:Npn \__lrnv_read_csv_row_to_seq:NN #1 #2
  {
    \ior_str_get:NN #1 \l__lrnv_tmp_str
    \seq_set_from_clist:NV #2 \l__lrnv_tmp_str
  }
\cs_generate_variant:Nn \seq_set_from_clist:Nn { NV }
\msg_new:nnn { lrnv } { file-not-found }
  { File~`#1'~not~found. }
\NewDocumentCommand \ReadCSVArray { m }
  {
    \file_if_exist:nTF {#1}
      { \__lrnv_csv_to_prop:n {#1} }
      { \msg_error:nnn { lrnv } { file-not-found } {#1} }
  }
\NewExpandableDocumentCommand \ArrayItem { m }
  { \prop_item:Nn \l__lrnv_array_prop { "#1" } }
\ExplSyntaxOff

\newcommand{\m}[1]{\mathsf{#1}}
\newcommand{\ttt}[1]{{\tt #1}}

\newcommand{\clang}{\ttt{C}}
\newcommand{\ccpp}{\clang/\cpp \ }




\newlist{protocol}{enumerate}{10}
\setlist[protocol]{label*={\arabic*.}, ref=\thesubsubsection.\arabic*}


\newcommand{\ifyestxt}{if {\textbf{yes}}}
\newcommand{\ifnotxt}{if {\color{BrickRed}\textbf{no}}}
\newcommand{\ifyes}{\internote{\ifyestxt}}
\newcommand{\ifno}{\internote{\ifnotxt}}
\newcommand{\gotosec}[1]{go to {\textbf{\nameref{#1}}}}
\newcommand{\nogo}[1]{\optgo{\ifnotxt}{#1}}
\newcommand{\yesgo}[1]{\optgo{\ifyestxt}{#1}}
\newcommand{\optgo}[2]{\internote{#1 then \gotosec{#2}}}



\newenvironment{protocolcond}[1]{#1\begin{addmargin}[1em]{0em} \begin{protocol}}{\end{protocol}\end{addmargin}}

\newcommand{\adjective}[1]{{\textit{#1}}}


\newcommand{\periodic}[1]{
\item \internote{select one} #1
    \begin{itemize}
    \item \textit{Daily}, \textit{Weekly}, \textit{Monthly}, \textit{Yearly}, \textit{Less than once a year}
    \end{itemize}
}
\newcommand{\ease}[1]{
\item \internote{select one} #1
    \begin{itemize}
    \item \textit{Extremely difficult}, \textit{Somewhat difficult}, \textit{Neither easy nor difficult}, \textit{Somewhat easy}, \textit{Extremely easy}
    \end{itemize}
}
\newcommand{\frequency}[1]{
\item \internote{select one} #1
    \begin{itemize}
    \item \textit{Never}, \textit{Sometimes}, \textit{About half the time}, \textit{Most of the time}, \textit{Always}
    \end{itemize}
}
\newcommand{\frequencyunsure}[1]{
\item \internote{select one} #1
    \begin{itemize}
    \item \textit{Unsure}, \textit{Never}, \textit{Sometimes}, \textit{About half the time}, \textit{Most of the time}, \textit{Always}
    \end{itemize}
}

\newcommand{\yesno}[1]{
\item \internote{select one} #1
    \begin{itemize}
    \item \textit{Definitely not}, \textit{Probably not}, \textit{Might or might not}, \textit{Probably yes}, \textit{Definitely yes}
    \end{itemize}
}

\newcommand{\responsessplit}[2]{
\vskip 0.5em
\begin{minipage}[t]{0.5\textwidth}
\begin{itemize}
\raggedright
#1
\end{itemize}
\end{minipage}\begin{minipage}[t]{0.5\textwidth}
\begin{itemize}
\raggedright
#2
\end{itemize}
\end{minipage}
\vskip 0.5em
}


\usepackage{changepage}
\usepackage{multicol}
\usepackage{multirow}
\usepackage{lscape}
\usepackage{flafter} 
\setlength{\multicolsep}{1.5pt plus .5pt minus .45pt}% 50% of original values
\newenvironment{responsestwocol}[1]
{
\begin{multicols}{2}
\itemize
}
{
\enditemize
\end{multicols}
}
\usepackage{float}

\usepackage{placeins}